% Generated draft: requires mathematical and source review.
\documentclass[11pt]{article}
\usepackage[margin=25mm]{geometry}
\usepackage{fontspec}
\setmainfont{Latin Modern Roman}
\usepackage{amsmath,amssymb,mathtools,unicode-math}
\setmathfont{Latin Modern Math}
\usepackage{enumitem,needspace,xurl,hyperref}
\setlength{\parindent}{0pt}
\setlength{\parskip}{5pt}
\setlength{\emergencystretch}{3em}
\newcommand{\N}{\mathbb{N}}
\newcommand{\Z}{\mathbb{Z}}
\newcommand{\Q}{\mathbb{Q}}
\newcommand{\R}{\mathbb{R}}
\newcommand{\C}{\mathbb{C}}
\newcommand{\F}{\mathbb{F}}
\newcommand{\A}{\mathbb{A}}
\newcommand{\PP}{\mathbb{P}}
\newcommand{\E}{\mathbb{E}}
\newcommand{\eps}{\varepsilon}
\newcommand{\dd}{\,\mathrm{d}}
\begin{document}
\section*{UnsolvedMath ID: 700348}
\textbf{Source problem code:} catalogue entry
\subsection{Definitions and mathematical statement}
Interpret the source notation \(C\) as the centralizer in the fixed ambient group \(G\); this interpretation is implicit rather than explicitly defined in the supplied record. For a subset \(S\subseteq G\), set
\[
C_G(S)=\{g\in G:gs=sg\text{ for every }s\in S\},\qquad C_G^2(S)=C_G(C_G(S)).
\]
For \(x\in G\), write \(C_G^2(x)=C_G^2(\{x\})\). A subgroup \(H\leq G\) is \emph{\(C\)-closed} if \(H=C_G^2(H)\), and is \emph{weakly \(C\)-closed} if
\[
\forall x\in H,\qquad C_G^2(x)\leq H.
\]
A group is \emph{locally finite} if every subgroup generated by finitely many elements is finite.

\textbf{Problem (V. A. Antonov and N. F. Sesekin, 1978).} Describe, up to group isomorphism, all locally finite groups \(G\) satisfying
\[
\forall H\leq G,\qquad H\neq G\ \Longrightarrow\ \bigl(\forall x\in H,\ C_G(C_G(\{x\}))\leq H\bigr).
\]
The quantifier ranges over all proper subgroups, including the trivial subgroup when it is proper, and all centralizers are taken in \(G\). Both finite and infinite locally finite groups are included.

The supplied background attributes to Gasch\"utz a study of finite groups whose proper subgroups are all \(C\)-closed. That background concerns the stronger closure condition and supplies no precise theorem or bibliographic reference for it.
\subsection{Short English statement}
Describe all locally finite groups in which every proper subgroup contains the ambient double centralizer of each of its elements. Here locally finite means that every finitely generated subgroup is finite. The centralizer meaning of \(C\) is implicit in the supplied source.
\subsection{Sources}
\begin{itemize}
\item Local UnsolvedMath registry, permanent ID \(700348\), exact source problem code \texttt{catalogue entry}, title \emph{Kourovka Notebook Problem 6.1}; supplied complete source record, upstream revision \texttt{d371ee917ac9647fbb06cd6fc529deec6f3dda82}.
\item \emph{The Kourovka Notebook (Unsolved Problems in Group Theory)}, original problem 6.1, sixth issue (1978); the supplied record identifies the twenty-first issue edition (2026), \texttt{arXiv:1401.0300v48}. Proposed by V. A. Antonov and N. F. Sesekin. The supplied collection description names E. I. Khukhro and V. D. Mazurov as editors.
\item The source background attributes the finite-group study for the \(C\)-closed condition to Gasch\"utz, without supplying a publication citation or a precise result.
\item The registry labels the problem \emph{open}. This is a supplied status label, not an independently verified literature assessment. No literature-verification date or explicit reference URL is supplied.
\end{itemize}
\end{document}
